\documentclass[../cours_1re.tex]{subfiles}
\begin{document}
\setcounter{chapter}{27}
\chapter{Complément : la géométrie des probabilités}
\label{ch:28}
\begin{objectif}
Complément : reconnaître la géométrie cachée dans les probabilités finies.
Prérequis : espérance, variance et projections du chapitre 27.
\end{objectif}
\section{Les variables aléatoires comme vecteurs}
Sur un univers fini dont toutes les issues ont un poids strictement positif,
les variables réelles forment un espace vectoriel. Posons
$\langle X,Y\rangle=\mathbb E(XY)$.
La bilinéarité et la symétrie sont immédiates ; $\mathbb E(X^2)>0$
pour $X\ne0$ puisque l'une de ses valeurs non nulles a un poids positif.
On a donc un produit scalaire.
Si certaines issues ont un poids nul, on les supprime, ou l'on identifie
les variables égales presque sûrement ; sinon la positivité stricte échouerait.

La projection de $X$ sur les constantes vaut $\mathbb E(X)$, car
$\langle X,1\rangle/\langle1,1\rangle=\mathbb E(X)$.
La variance est la norme au carré du résidu $X-\mathbb E(X)$.
Le théorème de meilleure constante du chapitre 22 est donc un théorème de projection.

\section{Prévoir à partir d'une information}
Une partition $A_1,\ldots,A_r$ de probabilités positives représente
l'information disponible : on sait dans quel groupe se trouve l'issue,
mais pas laquelle. Une prédiction utilisant seulement cette information
est une variable $Z$ constante, de valeur $z_i$, sur chaque $A_i$.
Définissons
\[
m_i=\frac{\mathbb E(X\mathbf1_{A_i})}{P(A_i)},\qquad
Y=\sum_{i=1}^r m_i\mathbf1_{A_i}.
\]
On appelle $Y$ l'espérance conditionnelle de $X$ sachant cette partition.
Elle est elle-même une variable aléatoire, alors que chaque $m_i$ est un nombre.

\begin{theoreme}[prédiction quadratique optimale]
Pour toute prédiction $Z$ constante sur les groupes,
\[
\mathbb E((X-Z)^2)=\mathbb E((X-Y)^2)+\mathbb E((Y-Z)^2).
\]
Le meilleur prédicteur au sens de l'erreur quadratique est donc $Y$,
unique à égalité presque sûre près.
\end{theoreme}
\begin{demonstration}
Par définition,
$\mathbb E((X-Y)\mathbf1_{A_i})
=\mathbb E(X\mathbf1_{A_i})-m_iP(A_i)=0$.
Le résidu est donc orthogonal à toute variable constante sur les groupes,
en particulier à $Y-Z$.
Développer $X-Z=(X-Y)+(Y-Z)$ et prendre l'espérance du carré :
le terme croisé s'annule. Le dernier carré est nul seulement si $Y=Z$
sur toutes les issues de poids positif.
\end{demonstration}

\section{Espérance et variance totales}
En sommant les égalités précédentes sur les groupes,
$\mathbb E(Y)=\mathbb E(X)$ : c'est la formule de l'espérance totale.
Définir la variance dans le groupe $i$ par
$v_i=\mathbb E((X-m_i)^2\mathbf1_{A_i})/P(A_i)$.
La décomposition orthogonale de $X-\mu=(X-Y)+(Y-\mu)$ donne
\[
\operatorname{Var}(X)=\sum_iP(A_i)v_i+\operatorname{Var}(Y).
\]
La dispersion totale se décompose en dispersion à l'intérieur des groupes
et dispersion entre leurs moyennes.

\begin{exemple}
Quatre issues équiprobables donnent $X=0,2,4,6$.
L'information distingue le groupe des deux premières et celui des deux dernières.
On a $Y=1,1,5,5$, $\mu=3$, $\operatorname{Var}(X)=5$.
Les variances internes valent chacune $1$, et $\operatorname{Var}(Y)=4$ :
$5=1+4$. Prédire par la moyenne globale $3$ donne une erreur quadratique $5$ ;
utiliser le groupe la réduit à $1$.
\end{exemple}

\section{Corrélation et dépendance}
Cauchy--Schwarz appliquée aux variables centrées donne
$|\operatorname{Cov}(X,Y)|\le\sigma(X)\sigma(Y)$.
Si les deux écarts types sont non nuls, le coefficient de corrélation
$\rho=\operatorname{Cov}(X,Y)/(\sigma(X)\sigma(Y))$ appartient à $[-1,1]$.
Il atteint $1$ ou $-1$ exactement lorsque les variables centrées sont
proportionnelles presque sûrement, par le cas d'égalité de Cauchy--Schwarz.
Une corrélation nulle exclut une liaison linéaire mesurée par ce produit,
mais n'exclut pas une dépendance non linéaire (exemple $Y=X^2$ du chapitre 22).
Elle ne démontre pas non plus une relation causale.

\section{Une première inégalité de Jensen}
Supposons $f$ dérivable sur un intervalle contenant les valeurs de $X$
et leur moyenne, avec $f'$ croissante. Le chapitre 11 montre que la courbe
est au-dessus de sa tangente en $\mu$ :
$f(x)\ge f(\mu)+f'(\mu)(x-\mu)$.
Prendre l'espérance donne
\[
\mathbb E(f(X))\ge f(\mathbb E(X)).
\]
C'est une forme de l'inégalité de Jensen pour une fonction convexe dérivable.
Avec $f(x)=x^2$, on retrouve $\mathbb E(X^2)\ge\mathbb E(X)^2$ ;
avec $f(x)=e^x$, $\mathbb E(e^X)\ge e^{\mathbb E(X)}$.
Pour prouver la convexité géométrique dans ce cadre, fixer $a<b$ et
$c=(1-t)a+tb$. Les pentes moyennes de $[a,c]$ et $[c,b]$ sont ordonnées
par les accroissements finis et la croissance de $f'$.
Après multiplication par les longueurs, cela donne
$f(c)\le(1-t)f(a)+tf(b)$, pour $0\le t\le1$.

\begin{finchapitre}
\section{Utilité et prolongements}
Le produit scalaire unifie les calculs de variance, les ajustements et
la prédiction avec information partielle.

En master, l'espérance conditionnelle sera définie pour des informations plus
générales que des partitions finies. Dans l'espace $L^2$, elle conservera
le sens de projection ; mesure, intégration et complétude seront alors nécessaires.
\begin{retenir}
Une information détermine les prédicteurs autorisés. L'espérance conditionnelle
est le meilleur d'entre eux pour l'erreur quadratique. Variance totale et Pythagore expriment la même structure.
\end{retenir}
\end{finchapitre}
\end{document}
